\Needspace{18\baselineskip}
\section{Cost relations and state refinement}

\begin{definition}[Cost preorder]\label{term:cost-preorder}
A cost order is a reflexive, transitive relation $Q$ on $\mathbb N$.
\end{definition}
\lean{Quantitative}{6}{9}{CostOrder}
The relation and its two preorder proofs are stored together. Antisymmetry,
decidability, and compatibility with addition are not part of this definition.

\begin{definition}[Reasonable triple]\label{term:reasonable}
A triple $(Q,P,S)$ of cost preorders is \emph{reasonable} if all four of the
following implications hold:
\begin{align*}
 a\,Q\,b\ \land\ c\,Q\,d &\ \Longrightarrow\ (a+c)\,Q\,(b+d),\tag{addition}\\
 a\,P\,b &\ \Longrightarrow\ b\leq a,\tag{decrease}\\
 a\,P\,b &\ \Longrightarrow\ a\,S\,b,\tag{inclusion}\\
 a\,P\,b\ \land\ b\,Q\,c\ \land\ c\,S\,d
   &\ \Longrightarrow\ a\,Q\,d.\tag{sandwich}
\end{align*}
\end{definition}
\lean{Quantitative}{12}{16}{Reasonable}
These four fields express the paper's three conditions: the decrease and
inclusion fields are the two parts of its condition~(B). The decrease field
will justify a strict descent in the soundness proof. Inclusion is used when
turning safety's auxiliary equivalences into an asymmetric up-to relation.

\Needspace{23\baselineskip}
\begin{definition}[Converse and the basic cost orders]\label{term:cost-orders}
The converse preorder is $a\,Q^{-1}\,b$ if $b\,Q\,a$. The four named
orders are
\[
 \begin{aligned}
 a\,E\,b &\iff a=b,&
 a\,I\,b &\iff b\leq a,\\
 a\,H\,b &\iff a\leq b,&
 a\,U\,b &\iff \text{true}.
 \end{aligned}
\]
Here $E$, $I$, $H$, and $U$ denote \texttt{exact}, \texttt{improvement},
\texttt{overhead}, and \texttt{any}, respectively.
\end{definition}
\lean{Quantitative}{20}{23}{CostOrder.converse}
\lean{Quantitative}{25}{25}{CostOrder.exact}
\lean{Quantitative}{27}{28}{CostOrder.improvement}
\lean{Quantitative}{30}{30}{CostOrder.overhead}
\lean{Quantitative}{32}{32}{CostOrder.any}
\begin{proof}[Verification of the preorder fields]
Converse preserves reflexivity. Given $b\,Q\,a$ and $c\,Q\,b$,
transitivity of $Q$ gives $c\,Q\,a$, establishing transitivity of
$Q^{-1}$. Equality is reflexive and transitive. The relation $I$ is
reflexive by $a\leq a$; from $b\leq a$ and $c\leq b$ one obtains
$c\leq a$. The corresponding verification for $H$ uses the usual order
direction. All requirements for the universal relation $U$ are immediate.
\end{proof}

\begin{lemma}[Exact costs are reasonable]
The triple $(E,E,E)$ is reasonable.
\end{lemma}
\lean{Quantitative}{34}{38}{CostOrder.reasonable_exact}
\begin{proof}
Substitute equals for equals in the addition condition. The decrease
condition reduces to $a\leq a$, inclusion is equality included in itself,
and the sandwich condition is transitivity of equality.
\end{proof}

\begin{lemma}[Improvement costs are reasonable]
The triple $(I,I,I)$ is reasonable.
\end{lemma}
\lean{Quantitative}{40}{44}{CostOrder.reasonable_improvement}
\begin{proof}
For addition, $b\leq a$ and $d\leq c$ give $b+d\leq a+c$. Decrease is
the definition of $I$, and inclusion is reflexive inclusion. The sandwich
premises give $d\leq c\leq b\leq a$, hence $a\,I\,d$.
\end{proof}

\begin{lemma}[Unconstrained termination costs are reasonable]
The triple $(U,I,U)$ is reasonable.
\end{lemma}
\lean{Quantitative}{46}{50}{CostOrder.reasonable_termination}
\begin{proof}
Addition, inclusion, and the sandwich condition have universally true
conclusions. Decrease is again the definition of $I$.
\end{proof}

\begin{lemma}[Overhead with exact auxiliary costs is reasonable]
The triple $(H,E,E)$ is reasonable.
\end{lemma}
\lean{Quantitative}{52}{56}{CostOrder.reasonable_overhead}
\begin{proof}
The addition condition is monotonicity of addition for $\leq$. Equality
implies the required weak decrease and includes into itself. In the sandwich
condition, substitute $a=b$ and $c=d$ into $b\leq c$ to obtain $a\leq d$.
\end{proof}

\begin{definition}[State refinement]\label{term:refinement}
For a cost preorder $Q$, define
\[
 s\preceq_Q t\quad\Longleftrightarrow\quad
 \begin{gathered}
 \text{for every $n,z$ with $s\steps{n}z$ and $\fin(z)$,}\\
 \text{there exist $m,w$ with $n\,Q\,m$, $t\steps{m}w$, and $\fin(w)$.}
 \end{gathered}
\]
State equivalence is mutual refinement:
$s\simeq_Q t$ if $s\preceq_Q t$ and $t\preceq_Q s$.
\end{definition}
\lean{Quantitative}{65}{67}{Machine.Refines}
\lean{Quantitative}{69}{70}{Machine.Equiv}
If $s$ has no successful execution, its refinement obligation is vacuous.
The definition therefore does not distinguish stuckness from divergence
unless a surrounding observation supplies additional distinctions.
For $Q=I$, the target must terminate in no more steps than the source.
For $Q=H$, it must terminate in at least as many steps. For $Q=E$, the
counts agree exactly; for $Q=U$, counts are ignored.

\begin{lemma}[Reflexivity]
For every $s$, $s\preceq_Q s$.
\end{lemma}
\lean{Quantitative}{74}{76}{Machine.Refines.refl}
\begin{proof}
Given a successful execution of length $n$ from $s$, choose that same
execution and final state as witnesses. Reflexivity supplies $n\,Q\,n$.
\end{proof}

\begin{lemma}[Transitivity]
If $s\preceq_Q t$ and $t\preceq_Q v$, then $s\preceq_Q v$.
\end{lemma}
\lean{Quantitative}{78}{83}{Machine.Refines.trans}
\begin{proof}
Take a successful execution of length $n$ from $s$. The first refinement
supplies a successful execution of length $m$ from $t$, with $n\,Q\,m$.
Apply the second refinement to this execution to obtain a successful
execution of length $k$ from $v$, with $m\,Q\,k$. Transitivity gives
$n\,Q\,k$, and these are the required witnesses. Determinism is unnecessary.
\end{proof}

\begin{lemma}[Monotonicity in the cost relation]
If $Q\subseteq P$ and $s\preceq_Q t$, then $s\preceq_P t$.
\end{lemma}
\lean{Quantitative}{85}{90}{Machine.Refines.mono}
\begin{proof}
Reuse each execution and final-state witness supplied by the original
refinement. Replace its cost proof $n\,Q\,m$ by $n\,P\,m$ using the inclusion.
\end{proof}
